% Physical pages: 430--436 (printed pages 407--413).
% Figures: none.
% Uncertainties: none.

$\prod_a(q_a-q'_a)$, so
\begin{equation}
\left(\prod_a(q_a-q'_a)\right)f(q)
=\left(\prod_a(q_a-q'_a)\right)f(q'),
\tag{9.5.40}\label{eq:9.5.40}
\end{equation}
which partly justifies our earlier remark that Eq. \eqref{eq:9.5.20} plays the role of a
delta function for integrals over the $q$s. Using Eq. \eqref{eq:9.5.32}, we now have
\[
\braket{q'}{f}=(-)^N\left[
 \int\left(\widetilde{\prod_b}dq_b\right)
 \left(\prod_a(q_a-q'_a)\right)f(q')\right].
\]
The term in the integrand proportional to $\prod_aq_a$ has coefficient
$f(q')$, so according to our definition of integration
$\braket{q'}{f}=(-)^Nf(q')$. Inserting this back in Eq. \eqref{eq:9.5.39} gives our
completeness relation
\[
\ket{f}=(-)^N\int\ket{q}
 \left(\widetilde{\prod_b}dq_b\right)\braket{q}{f},
\]
or as an operator equation
\begin{equation}
\mathbf{1}=\int\ket{q}
 \left(\widetilde{\prod_a}(-dq_a)\right)\bra{q}.
\tag{9.5.41}\label{eq:9.5.41}
\end{equation}
In exactly the same way, we can also show that
\begin{equation}
\mathbf{1}=\int\ket{p}
 \left(\widetilde{\prod_a}dp_a\right)\bra{p}.
\tag{9.5.42}\label{eq:9.5.42}
\end{equation}

We are now in a position to calculate transition matrix elements. As before,
we define time-dependent operators
\begin{equation}
Q_a(t)=\exp(iHt)Q_a\exp(-iHt)
\tag{9.5.43}\label{eq:9.5.43}
\end{equation}
\begin{equation}
P_a(t)=\exp(iHt)P_a\exp(-iHt)
\tag{9.5.44}\label{eq:9.5.44}
\end{equation}
and their right- and left-eigenstates
\begin{equation}
\ket{q;t}\equiv\exp(iHt)\ket{q},
\qquad
\ket{p;t}\equiv\exp(iHt)\ket{p},
\tag{9.5.45}\label{eq:9.5.45}
\end{equation}
\begin{equation}
\bra{q;t}\equiv\bra{q}\exp(-iHt),
\qquad
\bra{p;t}\equiv\bra{p}\exp(-iHt).
\tag{9.5.46}\label{eq:9.5.46}
\end{equation}
The scalar product between $q$-eigenstates defined at infinitesimally close
times is then
\[
\braket{q';\tau+d\tau}{q;\tau}
=\bra{q'}\exp(-iH\,d\tau)\ket{q}.
\]
Now insert Eq. \eqref{eq:9.5.42} to the left of the operator $\exp(-iH\,d\tau)$. It is
convenient here to define the Hamiltonian operator $H(P,Q)$ with all $P$s to
the left of all $Q$s, so that (for $d\tau$ infinitesimal)
\[
\bra{p}\exp[-iH(P,Q)d\tau]\ket{q}
=\braket{p}{q}\exp[-iH(p,q)d\tau].
\]
(We could move the c-number $H(p,q)$ to either side of the matrix element
without any sign changes because each term in the Hamiltonian is assumed to
contain an even number of fermionic operators.) This gives
\begin{align*}
\braket{q';\tau+d\tau}{q;\tau}
&=\int\braket{q'}{p}
 \left(\widetilde{\prod_a}dp_a\right)
 \bra{p}\exp(-iH\,d\tau)\ket{q}\\
&=\int\braket{q'}{p}
 \left(\widetilde{\prod_a}dp_a\right)\braket{p}{q}
 \exp[-iH(p,q)d\tau].
\end{align*}
Using Eqs. \eqref{eq:9.5.26} and \eqref{eq:9.5.27}, and noting that the products $p_aq_a$ and
$p_aq'_a$ commute with all anticommuting c-numbers, we find
\begin{equation}
\braket{q';\tau+d\tau}{q;\tau}
=\int\left(\widetilde{\prod_a}i\,dp_a\right)
 \exp\left[i\sum_ap_a(q'_a-q_a)-iH(p,q)d\tau\right].
\tag{9.5.47}\label{eq:9.5.47}
\end{equation}

The rest of the derivation follows the same lines as in Section 9.1. To
calculate the matrix element
$\bra{q';t'}\mathcal{O}_A(P(t_A),Q(t_A))
\mathcal{O}_B(P(t_B),Q(t_B))\cdots\ket{q;t}$ of a product of operators (with
$t'>t_A>t_B>\cdots>t$), divide the time-interval from $t$ to $t'$ into a large
number of very close time steps; at each time step insert the completeness
relation \eqref{eq:9.5.41}; use Eq. \eqref{eq:9.5.47} to evaluate the resulting matrix elements
(with $\mathcal{O}_A,\mathcal{O}_B$, etc. inserted where appropriate); move all
differentials to the left (this introduces no sign changes, because at each
step we have an equal number of $dp$s and $dq$s); and then introduce functions
$q_a(t)$ and $p_a(t)$ that interpolate between the values of $q_a$ and $p_a$ at
each step. We then find
\begin{align}
&\bra{q';t'}T\{\mathcal{O}_A(P(t_A),Q(t_A)),
 \mathcal{O}_B(P(t_B),Q(t_B)),\ldots\}\ket{q;t}
\notag\\
&\quad=(-i)^N\chi_N
 \int_{q_a(t)=q_a,\,q_a(t')=q'_a}
 \left(\widetilde{\prod_{a\tau}}dq_a(\tau)\,dp_a(\tau)\right)
 \notag\\
&\qquad\quad\cdot
 \mathcal{O}_A[p(t_A),q(t_A)]
 \mathcal{O}_B[p(t_B),q(t_B)]\cdots
 \notag\\
&\qquad\quad\cdot
 \exp\left[i\int_t^{t'}d\tau
 \left\{\sum_ap_a(\tau)\dot q_a(\tau)
 -H[p(\tau),q(\tau)]\right\}\right].
\tag{9.5.48}\label{eq:9.5.48}
\end{align}
The symbol $T$ here denotes the ordinary product if the times are in the order
originally assumed, $t_A>t_B>\cdots$. However, the right-hand side is totally
symmetric in the $\mathcal{O}_A,\mathcal{O}_B,\ldots$ (except for minus signs
where anticommuting c-numbers are interchanged) so this formula holds for
general times (between $t$ and $t'$), provided $T$ is interpreted as the
time-ordered product, with an overall minus sign if time-ordering the operators
involves an odd number of permutations of fermionic operators.

Up to this point we have kept track of the overall phase factor
$(-i)^N\chi_N$. But in fact these phases contribute only to the vacuum--vacuum
transition amplitude, and hence will not be of importance to us.

The transition to quantum field theory follows along the same lines as
described for bosonic fields in Section 9.2. The vacuum expectation value of a
time-ordered product of operators is given by a formula just like Eq. \eqref{eq:9.2.17}:
\begin{align}
&{}_{\mathrm{out}}\bra{\Omega}T\{\mathcal{O}_A[P(t_A),Q(t_A)],
\notag\\[-2pt]
&\qquad\qquad\mathcal{O}_B[P(t_B),Q(t_B)],\ldots\}
 \ket{\Omega}_{\mathrm{in}}
\notag\\
&\quad\propto\int
 \left[\prod_{\tau,\mathbf{x},m}dq_m(\mathbf{x},\tau)\right]
 \left[\prod_{\tau,\mathbf{x},m}dp_m(\mathbf{x},\tau)\right]
\notag\\
&\qquad\quad\cdot\mathcal{O}_A[p(t_A),q(t_A)]
 \mathcal{O}_B[p(t_B),q(t_B)]\cdots
\notag\\
&\qquad\quad\cdot\exp\left[i\int_{-\infty}^{\infty}d\tau
 \left\{\begin{aligned}
 &\int d^3x\sum_mp_m(\mathbf{x},\tau)\dot q_m(\mathbf{x},\tau)\\[-2pt]
 &\quad{}-H[q(\tau),p(\tau)]+i\epsilon\text{ terms}
 \end{aligned}\right\}\right],
\tag{9.5.49}\label{eq:9.5.49}
\end{align}
where the proportionality constant is the same for all operators
$\mathcal{O}_A,\mathcal{O}_B$, etc., and the `$i\epsilon$ terms' again arise
from the wave function of the vacuum. As before, we have replaced each
discrete index like $a$ with a space position $\mathbf{x}$ and a field index
$m$. We are also dropping the tilde on the product of differentials, since it
only affects the constant phase in the path integral.

A major difference between the fermionic and bosonic cases is that here we
will not want to integrate out the $p$s before the $q$s. Indeed, in the
standard model of electroweak interactions (and in other theories, such as the
older Fermi theory of beta decay) the canonical conjugates $p_m$ are auxiliary
fields unrelated to the $q_m$, and the Lagrangian is linear in the $\dot q_m$,
so that the quantity $\int d^3x\sum_mp_m\dot q_m-H$ in Eq. \eqref{eq:9.5.49} as it
stands is the Lagrangian $L$. Each term in the Hamiltonian for a fermionic
field that carries a non-vanishing quantum number (like the electron field in
quantum electrodynamics) generally contains an equal number of $p$s
(proportional to $q_m^\dagger$) and $q$s. In particular, the free-particle term
$H_0$ in the Hamiltonian is bilinear in $p$ and $q$, so that
\begin{align}
&\int_{-\infty}^{\infty}d\tau
 \left\{\int d^3x\sum_mp_m(\mathbf{x},\tau)\dot q_m(\mathbf{x},\tau)
 -H_0[q(\tau),p(\tau)]+i\epsilon\text{ terms}\right\}
\notag\\
&\qquad=-\sum_{mn}\int d^4x\,d^4y\,
 \mathcal{D}_{m\mathbf{x},n\mathbf{y}}p_m(x)q_n(y)
\tag{9.5.50}\label{eq:9.5.50}
\end{align}
with $\mathcal{D}$ some numerical `matrix'. The interaction Hamiltonian
$H_{\mathrm{int}}=H-H_0$ is a sum of products of equal numbers of fermionic
$q$s and $p$s (with coefficients that may depend on bosonic fields) so when we
expand Eq. \eqref{eq:9.5.49} in powers of the $H_{\mathrm{int}}$ we encounter a sum of
fermionic integrals of the form
\begin{align}
&\mathcal{I}_{n_1m_1n_2m_2\cdots n_Nm_N}
 (x_1,y_1,x_2,y_2,\ldots,x_N,y_N)
\notag\\
&\quad\equiv\int
 \left[\prod_{\tau,\mathbf{x},m}dq_m(\mathbf{x},\tau)\right]
 \left[\prod_{\tau,\mathbf{x},m}dp_m(\mathbf{x},\tau)\right]
 q_{m_1}(x_1)p_{n_1}(y_1)q_{m_2}(x_2)p_{n_2}(y_2)\cdots
\notag\\
&\qquad\quad\cdot q_{m_N}(x_N)p_{n_N}(y_N)
\notag\\
&\qquad\quad\cdot
 \exp\left(-i\sum_{mn}\int d^4x\,d^4y\,
 \mathcal{D}_{m\mathbf{x},n\mathbf{y}}p_m(x)q_n(y)\right),
\tag{9.5.51}\label{eq:9.5.51}
\end{align}
one such term for each possible set of vertices in the Feynman diagram, with
coefficients contributed by each vertex given by $i$ times the coefficient of
the product of fields in the corresponding term in the interaction.

To calculate this sort of integral, first consider a generating function for
all these integrals:
\begin{align}
\mathcal{J}(f,g)
={}&\int\left[\prod_{\mathbf{x},\tau,m}
 dq_m(\mathbf{x},\tau)\,dp_m(\mathbf{x},\tau)\right]
\notag\\
&\quad\cdot\exp\left\{\begin{aligned}
 &-i\sum_{mn}\int d^4x\,d^4y\,
 \mathcal{D}_{m\mathbf{x},n\mathbf{y}}p_m(x)q_n(y)\\[-2pt]
&-i\sum_m\int d^4x\,p_m(x)f_m(x)\\[-2pt]
&\qquad{}-i\sum_n\int d^4y\,g_n(y)q_n(y)
 \end{aligned}\right\},
\tag{9.5.52}\label{eq:9.5.52}
\end{align}
where $f_m(x)$ and $g_n(y)$ are arbitrary anticommuting c-number functions.
We shift to new variables of integration
\[
p'_m(x)=p_m(x)+\sum_n\int d^4y\,
 g_n(y)(\mathcal{D}^{-1})_{n\mathbf{y},m\mathbf{x}},
\]
\[
q'_n(y)=q_n(y)+\sum_m\int d^4x\,
 (\mathcal{D}^{-1})_{n\mathbf{y},m\mathbf{x}}f_m(x).
\]
Using the translation invariance condition \eqref{eq:9.5.36}, we then find
\begin{align}
\mathcal{J}(f,g)
={}&\exp\left(i\sum_{mn}\int d^4x\,d^4y\,
 (\mathcal{D}^{-1})_{n\mathbf{y},m\mathbf{x}}g_n(y)f_m(x)\right)
\notag\\
&\quad\cdot\int\left[\prod_{\mathbf{x},\tau,m}
 dq'_m(\mathbf{x},\tau)\,dp'_m(\mathbf{x},\tau)\right]
\notag\\
&\quad\cdot\exp\left(-i\sum_{mn}\int d^4x\,d^4y\,
 \mathcal{D}_{m\mathbf{x},n\mathbf{y}}p'_m(x)q'_n(y)\right).
\tag{9.5.53}\label{eq:9.5.53}
\end{align}
The integral is a constant (i.e., independent of the functions $f$ and $g$)
which can be shown using Eq. \eqref{eq:9.5.38} to be proportional to
$\operatorname{Det}\mathcal{D}$. Of more importance to us is the first factor.
Expanding this factor in powers of $gf$ and comparing with the direct expansion
of Eq. \eqref{eq:9.5.52}, we see that
\begin{align}
&\mathcal{I}_{n_1m_1n_2m_2\cdots n_Nm_N}
 (x_1,y_1,x_2,y_2,\ldots,x_N,y_N)
\notag\\
&\qquad\propto\sum_{\text{pairings}}\delta_{\text{pairing}}
 \prod_{\text{pairs}}
 \left(-i\mathcal{D}^{-1}\right)_{\text{paired }m\mathbf{x},n\mathbf{y}},
\tag{9.5.54}\label{eq:9.5.54}
\end{align}
with a proportionality constant that is independent of the $x$, $y$, $m$, or
$n$, and also independent of the number of these variables. The sum is over
all different ways of pairing $p$s with $q$s, not counting as different
pairings that only differ in the order of the pairs. In other words, we sum
over the $N!$ permutations either of the $p$s or the $q$s. The sign factor
$\delta_{\text{pairing}}$ is $+1$ if this permutation is even; $-1$ if it is
odd.

This sign factor and sum over pairings are just the same as we encountered in
our earlier derivation of the Feynman rules, with the sum over pairings
corresponding to the sum over ways of connecting the lines associated with
vertices in the Feynman diagrams, and the factors
$(\mathcal{D}^{-1})_{m\mathbf{x},n\mathbf{y}}$ playing the role of the
propagator for the pairing of $q_m(x)$ with $p_n(y)$. In the Dirac formalism
for spin $\frac12$, the free-particle action is
\begin{align}
&\int_{-\infty}^{\infty}d\tau
 \left\{\int d^3x\sum_mp_m(\mathbf{x},\tau)\dot q_m(\mathbf{x},\tau)
 -H_0[q(\tau),p(\tau)]\right\}
\notag\\
&\qquad=-\int d^4x\,\bar\psi(x)
 [\gamma^\mu\partial_\mu+m]\psi(x),
\tag{9.5.55}\label{eq:9.5.55}
\end{align}
where in the usual notation the canonical variables here are
\begin{equation}
q_m(x)=\psi_m(x),
\qquad
p_m(x)=-[\bar\psi(x)\gamma^0]_m=i\psi_m^\dagger(x)
\tag{9.5.56}\label{eq:9.5.56}
\end{equation}
with $m$ a four-valued Dirac index. Comparing this with Eq. \eqref{eq:9.5.50}, we find
here
\begin{align}
\mathcal{D}_{m\mathbf{x},n\mathbf{y}}
&=\left\{\gamma^0\left[\gamma^\mu\frac{\partial}{\partial x^\mu}
 +m-i\epsilon\right]\right\}_{mn}
 \delta^4(x-y)\notag\\
&=\int\frac{d^4k}{(2\pi)^4}
 \left\{\gamma^0[i\gamma^\mu k_\mu+m-i\epsilon]\right\}_{mn}
 e^{ik\cdot(x-y)}.
\tag{9.5.57}\label{eq:9.5.57}
\end{align}
(Though we shall not work it out in detail, the $i\epsilon$ term here arises
in much the same way as for the scalar field in Section 9.2.) The propagator is
then
\begin{equation}
(\mathcal{D}^{-1})_{m\mathbf{x},n\mathbf{y}}
=\int\frac{d^4k}{(2\pi)^4}
 \left\{[i\gamma^\mu k_\mu+m-i\epsilon]^{-1}[-\gamma^0]\right\}_{mn}
 e^{ik\cdot(x-y)},
\tag{9.5.58}\label{eq:9.5.58}
\end{equation}
just as we found in the operator formalism. The extra factor $-\gamma^0$
arises because this propagator is the vacuum expectation value of
\[
T\{\psi_m(x),-[\bar\psi(y)\gamma^0]_n\},
\qquad\text{not}\qquad
T\{\psi_m(x),\bar\psi_n(y)\}.
\]

As one example of a problem that is easier to solve by path-integral than by
operator methods, let us calculate the field dependence of the vacuum$\to$vacuum
amplitude for a Dirac field that interacts only with an external field. Take
the Lagrangian as
\begin{equation}
\mathcal{L}=-\bar\psi[\gamma^\mu\partial_\mu+m+\Gamma]\psi,
\tag{9.5.59}\label{eq:9.5.59}
\end{equation}
where $\Gamma(x)$ is an $x$-dependent matrix representing the interaction of
the fermion with the external field. According to Eq. \eqref{eq:9.5.49}, the vacuum
persistence amplitude in the presence of this external field is
\begin{align}
{}_{\mathrm{out}}\braket{\Omega}{\Omega}_{\mathrm{in},\Gamma}
\propto{}&\int
 \left[\prod_{\tau,\mathbf{x},m}dq_m(\mathbf{x},\tau)\right]
 \left[\prod_{\tau,\mathbf{x},m}dp_m(\mathbf{x},\tau)\right]
\notag\\
&\quad\cdot\exp\left\{-i\int d^4x\,
 p^T\gamma^0[\gamma^\mu\partial_\mu+m+\Gamma-i\epsilon]q\right\}
\tag{9.5.60}\label{eq:9.5.60}
\end{align}
with a proportionality constant that is independent of $\Gamma(x)$. We write
this as
\begin{align}
{}_{\mathrm{out}}\braket{\Omega}{\Omega}_{\mathrm{in},\Gamma}
\propto{}&\int
 \left[\prod_{\tau,\mathbf{x},m}dq_m(\mathbf{x},\tau)\right]
 \left[\prod_{\tau,\mathbf{x},m}dp_m(\mathbf{x},\tau)\right]
\notag\\
&\quad\cdot\exp\left\{-i\sum_{mn}\int d^4x\,d^4y\,
 p_m(x)q_n(y)\mathcal{K}[\Gamma]_{m\mathbf{x},n\mathbf{y}}\right\},
\tag{9.5.61}\label{eq:9.5.61}
\end{align}
where
\begin{equation}
\mathcal{K}[\Gamma]_{m\mathbf{x},n\mathbf{y}}
=\left\{\gamma^0[\gamma^\mu\partial/\partial x^\mu
 +m+\Gamma(x)-i\epsilon]\right\}_{mn}\delta^4(x-y).
\tag{9.5.62}\label{eq:9.5.62}
\end{equation}
To evaluate this, we change the variables of integration $q_m(x)$ to
\begin{equation}
q'_m(x)=\sum_n\int d^4y\,
 \mathcal{K}[\Gamma]_{m\mathbf{x},n\mathbf{y}}q_n(y).
\tag{9.5.63}\label{eq:9.5.63}
\end{equation}
The remaining integral is now $\Gamma$-independent, so the whole dependence of
the vacuum persistence amplitude is contained in the determinant arising
according to Eq. \eqref{eq:9.5.38} from the change of variables:
\begin{equation}
{}_{\mathrm{out}}\braket{\Omega}{\Omega}_{\mathrm{in},\Gamma}
\propto\operatorname{Det}\mathcal{K}[\Gamma].
\tag{9.5.64}\label{eq:9.5.64}
\end{equation}
To recover the results of perturbation theory, let us write
\begin{equation}
\mathcal{K}[\Gamma]\equiv\mathcal{D}+\mathcal{G}[\Gamma],
\tag{9.5.65}\label{eq:9.5.65}
\end{equation}
\begin{equation}
\mathcal{G}[\Gamma]_{m\mathbf{x},n\mathbf{y}}
=(\gamma^0\Gamma(x))_{mn}\delta^4(x-y),
\tag{9.5.66}\label{eq:9.5.66}
\end{equation}
and expand in powers of $\mathcal{G}[\Gamma]$. Eq. \eqref{eq:9.5.64} gives then
\begin{align}
{}_{\mathrm{out}}\braket{\Omega}{\Omega}_{\mathrm{in},\Gamma}
&\propto\operatorname{Det}
 \left(\mathcal{D}[\mathbf{1}+\mathcal{D}^{-1}\mathcal{G}[\Gamma]]\right)
\notag\\
&=[\operatorname{Det}\mathcal{D}]
 \exp\left(\sum_{n=1}^{\infty}\frac{(-1)^{n+1}}{n}
 \operatorname{Tr}(\mathcal{D}^{-1}\mathcal{G}[\Gamma])^n\right).
\tag{9.5.67}\label{eq:9.5.67}
\end{align}

This is just what we should expect from the Feynman rules: the contributions
from internal lines and vertices in this theory are $-i\mathcal{D}^{-1}$ and
$-i\mathcal{G}[\Gamma]$; the trace of the product of $n$ factors of
$-\mathcal{D}^{-1}\mathcal{G}[\Gamma]$ thus corresponds to a loop with $n$
vertices connected by $n$ internal lines; $1/n$ is the usual combinatoric
factor associated with such loops (see Section 6.1); the sign factor is
$(-1)^{n+1}$ rather than $(-1)^n$ because an extra minus sign is associated
with fermion loops; and the sum over $n$ appears as the argument of an
exponential because the vacuum persistence amplitude receives contributions
from graphs with any number of disconnected loops. The $\Gamma$-independent
factor $\operatorname{Det}\mathcal{D}$ is less easy to derive from the Feynman
rules; it represents the contribution of any number of fermion loops that
carry no vertices.

More to the point, a formula like Eq. \eqref{eq:9.5.64} allows us to derive
non-perturbative results by using topological theorems to derive information
about the eigenvalues of kernels like $\mathcal{K}[\Gamma]$. This will be
pursued further in Volume II.
